%% ma: L11-B10-0004
%% lop: 11
%% chuong: 4
%% bai: 10
%% dang: 11.10.5
%% loai: DS
%% muc: [NB, NB, TH, TH]
%% dap_an: "ĐĐSĐ"
%% nguon: "de-cuong-giua-ky-1-toan-11-nam-2026-2027-truong-thpt-hoang-van-thu-ha-noi-2-8.pdf (M08, Đề cương giữa kì 1 Toán 11 THPT Hoàng Văn Thụ, Hà Nội), Phần 2 Câu 7"
%% kien_thuc: "Bài 10 (giao tuyến của hai mặt phẳng)"
%% trang_thai: cho-duyet
%% ly_do: "Dạng mới đề xuất 11.10.5"
\begin{ex}
Cho hình chóp $S.ABCD$, biết $AB$ cắt $CD$ tại $E$, $AC$ cắt $BD$ tại $F$ trong mặt phẳng đáy. Khi đó:
\choiceTF{\True Đường thẳng $EF$ nằm trong mặt phẳng $(ABCD)$}{\True $AB$ là giao tuyến của hai mặt phẳng $(SAB)$ và $(ABCD)$}{$SF$ là giao tuyến của hai mặt phẳng $(SAB)$ và $(SCD)$, $SE$ là giao tuyến của hai mặt phẳng $(SAC)$ và $(SBD)$}{\True Gọi $G=EF\cap AD$, khi đó $SG$ là giao tuyến của mặt phẳng $(SEF)$ và mặt phẳng $(SAD)$}
\loigiai{a) Đúng: $E$, $F$ đều thuộc mặt phẳng đáy.

b) Đúng: $A$, $B$ là hai điểm chung của hai mặt phẳng.

c) Sai: $E\in AB\cap CD$ nên $(SAB)\cap(SCD)=SE$; $F\in AC\cap BD$ nên $(SAC)\cap(SBD)=SF$. Khẳng định đã đổi chỗ hai giao tuyến.

d) Đúng: $S$ và $G$ là hai điểm chung của $(SEF)$ và $(SAD)$.}
\end{ex}
